\documentclass[12pt]{article}
\usepackage{graphicx}
\usepackage{amsmath}
\usepackage{geometry}
\geometry{margin=1in}
\usepackage{hyperref}
\hypersetup{colorlinks=true, linkcolor=blue, urlcolor=blue}

\title{PET Image Uniformity Analysis Tool\\\vspace{0.5em}\large Program Documentation}
\author{Prepared by Ryan Byrd}
\date{\today}

\begin{document}
\maketitle

\section{Overview}
This program is a Python-based analysis tool designed to evaluate the axial uniformity of PET images using DICOM data exported from clinical PET/CT systems. It provides a quantitative measure of detector or reconstruction uniformity by calculating mean activity concentrations within a fixed circular region of interest (ROI) across all image slices. The results are expressed as percent deviation relative to the global mean over the central 85\% of the axial field of view (FOV). 

The tool was developed to align with the performance assessment principles outlined in NEMA NU~2 and ACR standards and manufacturer-specific PET uniformity tests.

\section{Program Workflow}
The program proceeds through the following logical steps:

\subsection{Step 1: DICOM Folder Selection}
The user is prompted to select a directory containing the reconstructed PET DICOM series. Each file in this directory represents a single transaxial slice of the PET volume. The program verifies that valid DICOM files are found and terminates otherwise.

\subsection{Step 2: DICOM Series Loading}
The program reads each slice using the \texttt{pydicom} library and sorts them according to their spatial position using the DICOM tag \texttt{ImagePositionPatient[2]} (Slice Number). The pixel spacing and rescale parameters (\texttt{RescaleSlope}, \texttt{RescaleIntercept}) are extracted to correctly convert pixel values to activity units (SUV or Bq/mL, depending on the dataset/manufacturer).

\subsection{Step 3: ROI Center Selection}
The middle slice of the dataset is displayed in a 250~mm~$\times$~250~mm cropped field of view centered on the image. The user is asked to click on the center of the uniform phantom region.

A live red crosshair cursor is overlaid to aid visual alignment. The clicked coordinates are stored as the ROI center for subsequent analysis.

\subsection{Step 4: ROI Definition}
A circular ROI of 180~mm diameter (i.e., radius 90~mm) is defined using the user-selected center. The radius is converted from millimeters to pixels based on the image’s \texttt{PixelSpacing} tag to ensure geometric accuracy.

\subsection{Step 5: ROI Mean Calculation}
For each slice in the series, the mean pixel value within the circular ROI is computed and stored. This produces an axial profile of mean ROI activity.

\subsection{Step 6: Central 85\% Slice Selection}
Only the central 85\% of slices are used to calculate the global mean, excluding peripheral slices that may contain edge effects or partial-volume artifacts. The global mean is then used to express each slice’s deviation as:
\[
\text{Deviation}(\%) = \left( \frac{M_i}{\bar{M}_{\text{central}}} - 1 \right) \times 100
\]
where $M_i$ is the mean activity in slice $i$, and $\bar{M}_{\text{central}}$ is the average over the central 85\%.

\subsection{Step 7: Output Directory Setup}
The user is prompted to select a base directory and provide a folder name for saving all output files. The program then automatically creates a results folder containing all data products.

\subsection{Step 8: CSV Export}
A CSV file (\texttt{PET\_ROI\_Results.csv}) is created, containing the slice number, mean ROI value, and percent deviation. This file can be imported into Excel or other statistical tools for further analysis.

\subsection{Step 9: Deviation Plot Generation}
A line plot of slice number vs.\ percent deviation is generated and saved as \texttt{PET\_DeviationPlot.png}. The plot includes shaded regions indicating:
\begin{itemize}
    \item $\pm5\%$: acceptable variation (green zone),
    \item $5\text{–}10\%$: warning zone (yellow),
    \item $>10\%$: failure region (red dashed lines).
\end{itemize}
The central 85\% region is highlighted in light blue to indicate the axial region considered for the purposes of this evaluation.

\subsection{Step 10: Splash View Generation}
A “splash view” is generated by displaying each slice (cropped to 250~mm~$\times$~250~mm) in a grid layout. The intensity is inverted (background white, structure black) for visual emphasis. The ROI circle is overlaid in red to show its position across slices. The image is saved as \texttt{PET\_SplashView.png}.

\section{Interpretation of Results}
The deviation plot provides a quantitative representation of PET system axial uniformity:
\begin{itemize}
    \item \textbf{Ideal uniformity:} All points lie within the $\pm5\%$ region.
    \item \textbf{Acceptable uniformity:} Deviations remain below $\pm10\%$.
    \item \textbf{Failure condition:} Any slice within the central 85\% exceeds $\pm10\%$.
\end{itemize}
Non-uniform deviations may indicate detector imbalance, improper normalization, attenuation correction errors, or issues with the reconstruction algorithm.

The splash view complements the quantitative results by allowing a qualitative assessment of uniformity across the entire volume.

\section{How to Use the Program}
\begin{enumerate}
    \item Launch the Python script \texttt{PET\_Image\_Uniformity.py}.
    \item When prompted, select the folder containing your PET DICOM files.
    \item A cropped 250~mm~$\times$~250~mm image will appear. Move your cursor over the uniform phantom region (the crosshair follows the mouse) and click once to define the ROI center.
    \item Choose a base folder where the results will be saved, and enter a name for the results subfolder.
    \item Wait for the analysis to complete. The terminal will display the output paths for:
    \begin{itemize}
        \item \texttt{PET\_ROI\_Results.csv} (numerical results)
        \item \texttt{PET\_DeviationPlot.png} (quantitative uniformity plot)
        \item \texttt{PET\_SplashView.png} (qualitative image grid)
    \end{itemize}
\end{enumerate}

\section{Notes and Recommendations}
\begin{itemize}
    \item Ensure that DICOM data represent a complete, reconstructed PET series.
    \item The tool assumes the ROI represents a uniform region of known activity (e.g., ACR PET phantom cylinder).
    \item Deviations beyond $\pm10\%$ should prompt recalibration or service verification of detector normalization.
    \item All pixel intensity units are linearly rescaled using DICOM header calibration factors; thus, results are robust to scanner-specific scaling.
\end{itemize}

\section{Conclusion}
This program automates the evaluation of PET image uniformity through an intuitive graphical workflow. It provides both quantitative metrics and visual confirmation to support physics testing, quality assurance, and service verification of PET systems.

\end{document}
